{\sloppy\footnotesize
\begin{longtable}{>{\raggedright\arraybackslash}p{0.062\textwidth}>{\raggedright\arraybackslash}p{0.07\textwidth}>{\raggedright\arraybackslash}p{0.408\textwidth}>{\raggedright\arraybackslash}p{0.2\textwidth}>{\raggedright\arraybackslash}p{0.1\textwidth}}
\caption{Every prediction, as its predictions file states it (shortened), with the expectation recorded before the runs and the verdict. Foreseen: a failure the predictions file recorded before the runs as expected to fail, as already fixed by earlier records or an energy balance, or as failing under a condition it named that then held. The verdicts of Stages 11 and 12 are their scorecards'.}\label{tab:s-predictions}\\
\toprule
Stage & & Prediction & Expectation & Verdict \\
\midrule
\endfirsthead
\toprule
Stage & & Prediction & Expectation & Verdict \\
\midrule
\endhead
1 & P1 & At 33.2 kJ and 500 mW/K, for both release durations, the hottest node reaches 0.15 K/min before the lumped node does, by more than half of the lumped time. & Central estimate: 4.2 s against 29 s for the 10 s release (86 \% earlier) and 11.5 s against 39 s for the 30 s release (71 \% earlier). & held \\
1 & P2 & The resolved mean temperature reaches 0.15 K/min in exactly 10 of 42 cases: 33.2 kJ at 100, 200 and 500 mW/K, 21.5 kJ at 200 and 500 mW/K, both releases. & none recorded & held \\
1 & P3 & The smallest busbar conductance at which the hottest node reaches 0.15 K/min while the lumped node does not is 100 mW/K (21.5 kJ, both release durations). & Central estimate: 100 mW/K. & held \\
1 & P4 & The hottest node reaches 10 K/min in at least one case in which the lumped node does not reach 1 K/min. & none recorded & held \\
2 & P5(a) & At 33.2 kJ and 500 mW/K, the reactive resolved cell's hottest region enters the internal-short band (255.7 C) while the reactive lumped cell does not. & P5(a) is therefore expected to fail on its second half (probability that the lumped cell stays out of the band: about one in three). & held \\
2 & P5(b) & At 33.2 kJ and 200 mW/K, the hottest region is pushed into the internal-short band (255.7 C) by its own reaction heat, while the reactive lumped cell stays out. & The lumped half is near certain; the resolved half is about even odds. & failed \\
2 & P5(c) & At 21.5 kJ, 500 mW/K and 30 s release, neither the hottest region nor the reactive lumped cell enters the internal-short band (255.7 C). & none recorded & held \\
3 & P1 & At 33.2 kJ, 30 s release and 500 mW/K, the hottest node reaches 0.15 K/min before the lumped node in more than 95 \% of kinetic draws. & Brief's central estimate: 99 \%. Study's central estimate: 100 \%. & held \\
3 & P2 & At 33.2 kJ, 30 s release, 500 mW/K, the hottest-node onset's central 95 \% interval is narrower than 20 s and the lumped node's wider than 20 s. & Brief's central estimate: 8 s and 40 s. Study's central estimate: 9.2 s and 64 s. & held \\
3 & P3 & The smallest busbar conductance at which only the local chemistry starts is 100 mW/K in more than half the kinetic draws, and never above 200 mW/K. & Brief's central estimate: 70 \% at 100 mW/K. Study's central estimate: 18.5 \% at 100 mW/K, and 10.5 \% at 500 mW/K: expected to fail on both counts. & failed, foreseen \\
3 & P4 & At 33.2 kJ, 30 s release, 500 mW/K, the hottest-node onset time's thermophysical standard deviation exceeds its kinetic one, ratio above 1.5. & Brief's central estimate: 3. Study's central estimate: none made. & failed \\
3 & P5 & At 33.2 kJ, 30 s release, 500 mW/K, the lumped onset time's kinetic standard deviation exceeds its thermophysical one, ratio above 1.5. & Brief's central estimate: 2. Study's central estimate: none made. & held \\
3 & P6 & Fewer than 10 \% of kinetic draws show the source paper's early solid electrolyte interphase outlier behaviour at the hottest node, an onset before the node reaches 100 C. & Brief's central estimate: below 10 \%. Study's central estimate: 36.5 \%: expected to fail. & failed, foreseen \\
4 & P1 & The anode-electrolyte submodel (holding the solid electrolyte interphase) carries over half the variance of the onset lead, lumped minus hottest-node onset, at 33.2 kJ, 30 s release, 500 mW/K. & Brief's central estimate: 65 \%. Study's expectation: Held, by a wide margin. & held \\
4 & P2 & The three submodels' main-effect shares of the onset lead's variance at 33.2 kJ, 30 s release, 500 mW/K sum to between 0.8 and 1.1 (interaction small). & Brief's central estimate: none given. Study's expectation: Uncertain, no number. & held \\
4 & P3 & With the anode-electrolyte submodel (holding the solid electrolyte interphase) at its posterior mean, the share of draws whose hottest node starts below 100 C falls below 5 \%. & Brief's central estimate: none given. Study's expectation: Held. & held \\
4 & P4 & Restricting the dominant submodel to the half of its posterior samples nearest the mean narrows the onset lead's 2.5 to 97.5 \% width by more than 40 \%. & Brief's central estimate: none given. Study's expectation: May fail. & failed \\
4 & P5 & The ordering of the three submodels by their share of the onset lead's variance is the same at 100 mW/K as at 500 mW/K. & Brief's central estimate: none given. Study's expectation: Held, if the anode-electrolyte file leads at 500 mW/K. & held \\
5 & P1 & At the central contact conductance with the posterior-mean kinetics, the resolved neighbour propagates in series and parallel, its time to propagation within the measured 147 to 342 s. & Study's expectation: May fail: by the estimates above the hot spot stays below 255.7 C, or, if the line of contact rises further, reaches it within the first minute. & failed \\
5 & P2 & In both connections, series and parallel, the lumped neighbour propagates more than 20 s later than the resolved one, or not at all within 900 s. & Study's expectation: Held, most likely as ``not at all'': its passive peak should lie near 145 to 170 C. & held \\
5 & P3 & At the resolved model's propagation time, its mid-surface temperature lies between 150 and 255.7 C: the hot spot leads the surface a thermocouple sees. & Study's expectation: Toward the trigger, held if propagation occurs (that face is the line of contact); away from it, may fail below 150 C. & failed \\
5 & P4 & The resolved model's mean mid-surface heating rate over the 60 s before propagation lies between 0.2 and 0.8 C/s (measured about 0.4). & Study's expectation: Uncertain. & failed \\
5 & P5 & The predicted series and parallel times to propagation differ by less than 30 s (the measured means differ by 8 s). & Study's expectation: Held if both propagate: the tabs carry 4 to 8 \% of the contact path. & failed \\
5 & P6 & The share of the 4,800 factorial runs in which the resolved neighbour propagates lies between 30 and 90 \% (measured 9 of 16, 56 \%). & Study's expectation: May fail low: the model lacks the vent, ejecta and flame heat, and the estimates above fall short of 255.7 C at the central contact. & failed \\
5 & P7 & The contact and tab conductance range carries more of the spread in predicted time to propagation than the kinetic posterior (standard deviation ratio above 1), reversing Stage 3's onset finding. & Study's expectation: Held if enough runs propagate to give both spreads: the contact range changes the path 4.6-fold. & failed \\
6 & P1 & In Stage 5 runs the lumped neighbour first reaches 150 C within the measured 147 to 342 s in over half its crossings; the hot spot always before 147 s. & Brief's expectation: Likely; the Stage 5 central case already shows it. Study's expectation: Fails on its first part. & failed, foreseen \\
6 & P2 & The face away from the trigger, at the thermocouple's position, exceeds 133 C, the largest neighbour maximum of five non-propagating tests, in fewer than half the Stage 5 runs. & Brief's expectation: Uncertain. Study's expectation: Fails. & failed, foreseen \\
6 & P3 & With the trigger duration matched to the measured trigger-to-peak times, the lumped neighbour's peak rises above its Stage 5 central value (159 to 162 C) by at least 20 K. & Brief's expectation: Likely; a longer release keeps the neighbour heated past its Stage 5 peak time. Study's expectation: Fails. & failed, foreseen \\
6 & P4 & With the trigger energy raised to the Battery Failure Databank's total including the ejected fraction, the model's trigger peak falls inside the measured range of the sixteen trigger peaks. & Brief's expectation: Uncertain; the Databank's split and the tests' cells need not agree. Study's expectation: Fails. & failed, foreseen \\
6 & P5 & In the case matched to the tests in duration and energy, the resolved hot spot reaches 255.7 C in at least one run, and the lumped neighbour in none. & Brief's expectation: Close to even on the first part; likely on the second. Study's expectation: Fails on its first part; holds on its second. & failed, foreseen \\
6 & P6 & In the case matched to the tests, over the 200 kinetic posterior draws, the hot spot reaches 255.7 C in between 10 and 70 \% of them (central 35 \%). & Brief's expectation: Uncertain. Study's expectation: Fails, at 0 \%. & failed, foreseen \\
6 & P7 & In no case does the lumped neighbour reach 255.7 C before the hot spot does. & Brief's expectation: Likely. Study's expectation: Holds: where both reach 255.7 C (at the total energy, if anywhere), the line of contact leads. & held \\
7 & P1 & At the central contact with per-test triggers, C1 (hot spot at 255.7 C) predicts propagation in fewer than 9 of the 16 tests. & Brief's expectation: Likely; Stage 6 ignited the hot spot only at the Databank's total energy and, at the median peak, only at the high contact. Study's expectation: Holds, with eight. & held \\
7 & P2 & At the central contact, C3 (lumped neighbour at 150 C) predicts propagation in 9 or more of the 16 tests and gets at least 10 of 16 outcomes right. & Brief's expectation: Uncertain; the paper's own criterion on the paper's own trigger peaks. Study's expectation: Fails on its second part, by one test. & failed, foreseen \\
7 & P3 & Among the propagated tests that C3 predicts, the predicted crossing time lies within 60 s of the measured propagation time in at least half. & Brief's expectation: Uncertain. Study's expectation: Fails. & failed, foreseen \\
7 & P4 & C4 (separator melting at the hot spot, 138.5 C) predicts propagation in more tests than were measured (more than 9 of 16) at the central contact. & Brief's expectation: Likely; the hot spot passed 138.5 C in every Stage 5 run at the central and high contacts. Study's expectation: Holds. & held \\
7 & P5 & An aged solid electrolyte interphase lowers the hot spot's first self-heating temperature but changes no C4 or C5 verdict at the central contact. & Brief's expectation: Likely, conditional on A4. Study's expectation: Fails on its first part; its second part is expected to hold. & failed, foreseen \\
7 & P6 & Ranking the sixteen tests by trigger energy ranks the neighbour's lumped peak in the same order within each connection (no inversions). & Brief's expectation: Likely; a check on the mapping, not a result. Study's expectation: Holds. & held \\
7 & P7 & The series and parallel halves receive the same predicted propagation count within one, for every criterion, at the central contact. & Brief's expectation: Likely; the model distinguishes the connections by a tab only, and the measured counts (5 of 8 and 4 of 8) differ by one. Study's expectation: Fails, on C1 (and C6). & failed, foreseen \\
7 & P8 & In the internal-short event runs at the central contact, the resolved neighbour's peak after the event lies within the propagating neighbours' measured 381 to 763 C in at least half. & Brief's expectation: Uncertain. Study's expectation: Fails, with no run inside. & failed, foreseen \\
8 & P1 & Of the nine propagated tests, f*(C1) is 0 in the six Stage 7 hits and below 0.3 in the other three. & Brief's expectation: Likely; the hot spot is close to 255.7 C in several misses. Study's expectation: Holds. & held \\
8 & P2 & Of the seven non-propagated tests, at least three, including the two Stage 7 false alarms at 0, have f*(C1) below 0.3, so the propagated and non-propagated tests' f*(C1) ranges overlap. & Brief's expectation: Likely. Study's expectation: Holds, with all seven below 0.3. & held \\
8 & P3 & f*(C2) lies between 0.05 and 0.5 in at least twelve of the sixteen tests. & Brief's expectation: Uncertain. Study's expectation: Fails. & failed, foreseen \\
8 & P4 & With 0.2 of the ejected energy delivered, C2 predicts propagation in at least nine tests, crossing within 60 s of the measured time in at least four of its hits. & Brief's expectation: Uncertain; Stage 7 had one of eight within 60 s for C3 at f = 0. Study's expectation: Fails on its second part. & failed, foreseen \\
8 & P5 & At each common fraction (0.1, 0.2 and 0.3), no criterion scores more than 12 of 16: delivered heat moves the predicted count but not past the cap the outcomes set. & Brief's expectation: Likely. Study's expectation: Holds. & held \\
8 & P6 & Within each connection, f*(C1) and f*(C2) never increase with trigger energy (no inversion). & Brief's expectation: Likely; a check on the mapping. Study's expectation: Holds, with ties. & held \\
8 & P7 & Joule heat of 0.2 of the stored electrical energy (the electrical bracket) lowers f*(C1) of the parallel tests by less than 0.1 in every test. & Brief's expectation: Uncertain. Study's expectation: Fails, on M50\_PT\_Exp1. & failed, foreseen \\
8 & P8 & Delivering through the quarter of the facing half nearest the cap lowers f*(C1) in every sensitivity test where f* is above 0 and raises f*(C2) in none. & Brief's expectation: Likely for the first part, uncertain for the second. Study's expectation: Holds. & held \\
9 & P1 & With the 21.39 s delivery, C1 reproduces no propagated test's measured time within 30 s at any fraction at or above f*'s lower end. & Brief's expectation: Likely (the hot spot crosses at 11.9 to 18.5 s at the threshold's upper end). Study's expectation: Holds. & held \\
9 & P2 & With the 21.39 s delivery, C2 reproduces at least six of nine propagation times within 30 s, at a fraction between f*'s lower end and 0.1 above its upper end. & Brief's expectation: Uncertain (seven upper-end crossings lie inside the window; ST\_Exp3's lies at 61.4 s). Study's expectation: Holds. & held \\
9 & P3 & Every C2 time-matching fraction found in P2 lies no more than 0.02 above the upper end of f*(C2). & Brief's expectation: Uncertain. Study's expectation: Uncertain; it turns on M50\_PP\_Exp1. & held \\
9 & P4 & For the five low-maximum non-propagated tests, the peak-matching fraction at the thermocouple node lies below both f*(C1) and f*(C2). & Brief's expectation: Likely. Study's expectation: Fails, on M50\_ST\_Exp1 and ST\_Exp2, whatever the runs show. & failed, foreseen \\
9 & P5 & The peak-matching fractions (five non-propagated tests) and the C2 time-matching fractions (nine propagated tests) separate: the largest of the former is below the smallest of the latter. & Brief's expectation: Uncertain; this is the stage's question. Study's expectation: Fails. & failed, foreseen \\
9 & P6 & With a 240 s constant-power delivery, the C1 time-matching fraction (hot spot at 255.7 C at the measured time) exceeds 0.2 in M50\_PT\_Exp4, M50\_PP\_Exp2 and M50\_ST\_Exp3. & Brief's expectation: Uncertain. Study's expectation: Uncertain; fails if M50\_PT\_Exp4 has no f\_t. & failed, foreseen \\
9 & P7 & The bracket on the ejected-to-body energy ratio rescales every time- and peak-matching fraction by 0.75 to 1.40 and changes P5's verdict in neither direction. & Brief's expectation: Likely. Study's expectation: Holds. & held \\
9 & P8 & An aged solid electrolyte interphase changes the C2 time-matching fraction by less than 0.01 in M50\_PP\_Exp4, M50\_SP\_Exp1 and M50\_SP\_Exp2. & Brief's expectation: Likely (Stage 7: the aged SEI changed no verdict). Study's expectation: Holds. & held \\
10 & P1 & Before propagation, the model under the 21.39 s delivery, at one end of its threshold bracket, fits M50\_PT\_Exp4's measured history at the Stage 9 (far-face) node within 25 K root-mean-square. & Brief's expectation: uncertain. & held \\
10 & P2 & The long-delivery readings (240 and 360 s) fit M50\_PT\_Exp4's measured history with a smaller root-mean-square difference than the short-delivery reading. & Brief's expectation: uncertain; the measured history's shape decides, and A0.2 is recorded before this prediction is scored, not before it is written. & failed \\
10 & P3 & M50\_PT\_Exp4's measured history rises over 30 K in the 60 s before propagation, which the short delivery reproduces within a factor of two and the 360 s delivery does not. & Brief's expectation: uncertain. & failed, foreseen \\
10 & P4 & The peak-matching fraction is 0 at all three nodes in the same four low-maximum tests as at the Stage 9 node; M50\_PT\_Exp1's facing-node value is below its Stage 9 value. & Brief's expectation: likely (the far face is the coolest of the three). & held \\
10 & P5 & The separation fails at every node: the largest peak-matching fraction (five non-propagated tests) exceeds the smallest C2 time-matching fraction at the Stage 9, 90 degree and facing nodes. & Brief's expectation: likely. & failed, foreseen \\
10 & P6 & Held out, no predictor beats one probability for every test by more than 0.1 in mean log score in either split direction, and no permutation p-value is below 0.05. & Brief's expectation: likely. & held \\
10 & P7 & No ranked model quantity correlates with the measured outcome or time at p below 0.05 by permutation test, and the connection splits none of them. & Brief's expectation: likely. & held \\
10 & P8 & The survey finds at least one dataset with per-test outcomes and times for cells in direct contact, and none whose cell this study's parameters cover without an invented value. & Brief's expectation: uncertain. & failed \\
11 & P1 & For each of the four tests, f*(C2) under every kinetic draw lies within a factor of two of its posterior-mean value, or within 0.025 of it where that is 0. & Brief's expectation: likely. & failed \\
11 & P2 & For each of the four tests, f*(C1) under every kinetic draw lies within a factor of two of its posterior-mean value, or within 0.025 of it where that is 0. & Brief's expectation: uncertain; the hot spot's onset is local and the kinetics set it. & held \\
11 & P3 & Under each of the eight draws, the range of the C2 threshold energy across the four tests exceeds the largest within-test band of that energy across the draws. & Brief's expectation: likely. & held \\
11 & P4 & In the four tests with zero peak-matching fraction, the far face with no delivery exceeds the measured maximum under at least three of four thermal draws and at both contacts. & Brief's expectation: uncertain; the low contact and the smallest radial conductivity are the two that could remove it. & failed, foreseen \\
11 & P5 & In each of the four tests the band of f*(C2) across the four thermal draws is narrower than its band across the four kinetic draws. & Brief's expectation: uncertain. & failed \\
11 & P6 & In at least two of six tests the posterior mean matched within 10 s, a kinetic draw at that fraction removes the crossing or shifts it over 60 s. & Brief's expectation: likely. & failed \\
12 & P1 & With E\_b capped at 33.2 kJ, conduction and radiation alone meet C1 in none of the eight tests at the central contact. & Brief's expectation: likely (Stage 5 found no ignition at 33.2 kJ with the two-cell study's trigger); the file records it as fixed by Stage 6's records. & held \\
12 & P2 & At the high contact the same holds. & Brief's expectation: uncertain; the file records it as close to fixed (Stage 6, with other tabs). & held \\
12 & P3 & With E\_e fixed at the Databank's median, the delivered energies E* across the nine propagated tests still span at least tenfold. & Brief's expectation: likely; the file records it as fixed by Stage 9's records. & held \\
12 & P4 & With E\_e fixed, the propagated and non-propagated tests' fractions still overlap (the largest f\_max above the smallest f\_t(C2)). & Brief's expectation: likely; the file records it as fixed by Stage 9's records. & held \\
\bottomrule
\end{longtable}
}
